
While the Agent-Based Model (ABM) served as the simulation environment, \textbf{Deep Reinforcement Learning (DRL)} was the critical methodology required to autonomously discover optimal, budget-constrained policies within that complex system \cite{ZhengAI2022}. DRL integrated the decision-making framework of Reinforcement Learning with the representation learning capabilities of Deep Neural Networks (DNNs).

In this setup, the Municipal LGU agent utilized a neural network to approximate the optimal policy, learning to map high-dimensional state inputs—such as the heterogeneous psychological states of thousands of households—to precise adaptive decisions \cite{Hertweck2023Values}. This integration was essential because traditional tabular RL methods could not scale to the massive, continuous state spaces generated by complex socio-environmental models \cite{Mousavi2021Overview}.

\subsection{The ABM as a Markov Decision Process}
\label{subsec:abm_mdp}

To enable autonomous optimization, the literature supported framing the ABM simulation not merely as a model, but as a stochastic environment formally defined as a \textbf{Markov Decision Process (MDP)}. In this hybrid architecture, the ABM functioned as a high-fidelity data generator, providing the state representation ($S_t$)—encapsulating household compliance levels and budget statuses—and processing the policy agent’s actions ($A_t$) to generate the next state ($S_{t+1}$) and a corresponding reward signal ($R_t$). Studies in urban resource management \cite{Rajesh2025} and smart city logistics \cite{HaMinh2025} demonstrated that this approach provided the ``experience replay'' data required to train Deep RL agents in the absence of pre-existing datasets \cite{Kompella2017}.

\subsection{The Sparse Reward Problem in Policy Optimization}
\label{subsec:sparse_reward}

While DRL demonstrated superhuman performance in complex control environments \cite{Silver2016}, its application to public policy and resource allocation presented unique computational challenges. Unlike game environments with immediate feedback scores, governance models often suffered from delayed causal effects, where a policy implemented today might not have yielded observable compliance changes for months.

This phenomenon created a \textbf{Sparse Reward Environment}, which necessitated advanced exploration strategies beyond standard random-walk algorithms. The ``Sparse Reward Problem'' referred to the difficulty RL agents faced when feedback signals were infrequent or delayed \cite{SuttonBarto2018}. In high-dimensional search spaces—such as a municipal budget allocation across multiple barangays ($d=21$)—a standard agent acted as a ``random wanderer,'' unlikely to stumble upon the precise sequence of actions required to trigger a positive outcome. Research by \cite{Andrychowicz2017} presented at NeurIPS demonstrated that in such environments, standard agents often failed to learn entirely, converging to a ``lazy'' local optimum (e.g., doing nothing/Status Quo) because the probability of randomly discovering a successful strategy was near zero.

\subsection{Heuristic-Guided Exploration}
\label{subsec:hurl}

To overcome the inefficiency of random exploration in sparse environments, recent literature advocated for \textbf{Heuristic-Guided Reinforcement Learning (HuRL)}. \cite{Cheng2021} defined HuRL as the integration of domain-specific rules to prune the action space, allowing the agent to focus on ``promising'' regions of the solution landscape.

Crucially, HuRL did not replace the learning process; it accelerated it. As noted by \cite{Vecerik2018} in their work published in AAAI, heuristics acted as ``saliency filters'' that directed the agent's attention to relevant state-variables (e.g., non-compliant areas) while leaving the optimization of \textit{intensity} and \textit{duration} to the neural network. This hybrid approach—combining human domain knowledge with machine optimization—was proven to outperform pure ``tabula rasa'' learning in resource-constrained environments \cite{Nair2018}. 

\subsection{Reward Shaping}
\label{subsec:reward_shaping}

A third critical component in modern DRL architectures was \textbf{Reward Shaping}, formally introduced by \cite{Ng1999}. Reward shaping involved augmenting the raw environmental signal with intermediate rewards to guide the agent toward desirable behaviors.

However, \cite{Amodei2016} warned of ``Reward Hacking,'' where an agent optimized the shaped reward rather than the true objective. To mitigate this, \cite{Popov2017} suggested the use of ``Potential-Based Reward Shaping,'' where shaped rewards were strictly tied to physics-based thresholds. In the context of behavioral change, this aligned with the ``Threshold Hypothesis'' in sociology \cite{Granovetter1978}, which suggested that interventions had to exceed a critical ``dose'' to trigger collective action. By encoding this threshold into the reward function, DRL agents could be trained to avoid the ``Dilution Error''—the common administrative failure of spreading resources too thinly to have an effect.

\subsection{Current Trends and Research Gap}
\label{subsec:research_gap}

Current applications of DRL in waste management remained predominantly focused on technical, industrial, or hardware-based optimization, rather than policy governance. For instance, recent advancements extensively demonstrated the efficacy of DRL in automating waste classification and detection. Duhayyim et al. \cite{Duhayyim2022} and Khan et al. \cite{Khan2024} utilized Deep Q-Networks (DQN) and Mask R-CNN to automate the complex visual task of waste object detection, achieving high accuracy in segregating recyclables at the processing stage. Similarly, Udayakumar et al. \cite{Udayakumar2023} integrated Improved Particle Swarm Optimization (IPSO) with MobileNetV2 to enhance the precision of waste categorization in smart city frameworks.

Beyond classification, DRL was applied to control industrial processes, such as waste biorefining \cite{Gao2024} and plant machinery optimization \cite{Kumar2022}. In logistics, stochastic optimization was used for vehicle routing \cite{Akbarpour2021, Khallaf2024} and reverse logistics \cite{Karagoz2022}. This review highlighted a significant gap: while DRL was widely used for \textit{physical} waste sorting, its application for \textit{policy} sorting—optimizing the allocation of governance resources to influence human behavior—remained under-explored.
